\documentclass[11pt,a4paper]{article}
\usepackage[utf8]{inputenc}
\usepackage{amsmath}
\usepackage{amsfonts}
\usepackage{amssymb}
\usepackage{graphicx}
\usepackage{geometry}
\usepackage{hyperref}
\usepackage{enumitem}
\usepackage{booktabs}
\usepackage{xcolor}

\geometry{margin=1in}

\title{\textbf{FVCOM-MP Microplastic Infiltration Physics Model}\\
\large Mathematical Formulation and Parameterization}
\author{WaterPACT Development Team}
\date{\today}

\begin{document}

\maketitle

\tableofcontents
\newpage

\section{Overview}

This document presents the complete mathematical formulation for implementing 1D infiltration of microplastics (MP) and associated toxics into sediment columns within the FVCOM-MP (Finite Volume Community Ocean Model - MicroPlastic) framework. The model incorporates sophisticated retention mechanisms based on DLVO theory and strain-induced blockage effects.

\section{Seepage Flow Input}

The model requires vertical seepage velocity (Darcy velocity) input with two implementation options:

\subsection{Implementation Options}
\begin{enumerate}
    \item \textbf{Uniform Option}: Single $w_D$ value applied to entire computational domain via namelist parameter
    \item \textbf{Spatial Varying Option}: Node-specific $w_D$ values read from input file for spatial heterogeneity
\end{enumerate}

\section{Governing Equations for Microplastic Infiltration}

The infiltration model employs a two-phase approach distinguishing between movable and blocked microplastics within the sediment matrix.

\subsection{Primary MP Transport Equation}

The transport equation for movable microplastics ($p_{MP}$) incorporates advection, diffusion, blocking, and leaching processes:

\begin{equation}
\boxed{
\frac{\partial p_{MP}}{\partial t} + (w_D + w_{ss}) \frac{\partial p_{MP}}{\partial z} = \frac{\partial}{\partial z}\left[D_L \frac{\partial p_{MP}}{\partial z}\right] - \lambda|w_D + w_{ss}|p_{MP} - S_{leach}^p
}
\end{equation}

\subsection{Blocked MP Equation}

The evolution equation for blocked microplastics ($s_{MP}$) accounts for particle retention and release:

\begin{equation}
\boxed{
\frac{\partial s_{MP}}{\partial t} = \lambda|w_D + w_{ss}|p_{MP} - S_{leach}^s
}
\end{equation}

\subsection{Total Concentration}

The total microplastic concentration in the sediment is:

\begin{equation}
\boxed{
c_{MP} = p_{MP} + s_{MP}
}
\end{equation}

\subsection{Variable Definitions}

\begin{itemize}
    \item $p_{MP}$: Movable MP concentration in sediment [mass/volume]
    \item $s_{MP}$: Blocked MP concentration (immobilized by solid matrix) [mass/volume]
    \item $c_{MP}$: Total MP concentration in sediment [mass/volume]
    \item $w_D$: Darcy velocity (seepage flow, positive/negative) [m/s]
    \item $w_{ss}$: MP settling velocity within solid matrix [m/s]
    \item $z$: Vertical coordinate (sigma coordinate system) [m]
    \item $t$: Time [s]
    \item $\lambda$: Effective percolation blockage rate [1/s]
    \item $D_L$: Mechanical seepage diffusivity [m$^2$/s]
    \item $S_{leach}^p$, $S_{leach}^s$: Leaching source/sink terms [mass/(volume·time)]
\end{itemize}

\section{Mechanical Diffusivity Model}

The mechanical seepage diffusivity combines mechanical dispersion and molecular diffusion effects:

\begin{equation}
\boxed{
D_L = c \cdot d_g |w_D + w_{ss}| + \frac{D_m}{\tau}
}
\end{equation}

\subsection{Parameters}
\begin{itemize}
    \item $c$: Solid-matrix mechanical diffusivity coefficient [dimensionless]
    \item $d_g$: Solid-matrix grain size [m]
    \item $D_m$: Molecular diffusivity [m$^2$/s]
    \item $\tau$: Tortuosity in porous media [dimensionless]
\end{itemize}

\section{Percolation Blockage Rate Model}

The total blockage rate combines DLVO theory and strain-induced effects:

\begin{equation}
\boxed{
\lambda = \lambda_{DLVO} + \lambda_{strain}
}
\end{equation}

\subsection{DLVO Theory Component}

Based on colloid filtration theory incorporating van der Waals forces and electrostatic interactions:

\begin{equation}
\boxed{
\lambda_{DLVO} = \frac{3(1-\phi)}{2d_g}[\eta_D + \eta_I + \eta_G]\alpha_{primary}
}
\end{equation}

where $\alpha_{primary} = 1$ when Hamaker constant $A \geq 0$, otherwise $\alpha_{primary} = 0$.

\subsubsection{Retention Efficiency Components}

\paragraph{1. Diffusion-induced Retention}
Accounts for van der Waals attraction forces and Brownian motion:
\begin{equation}
\boxed{
\eta_D = 2.4 A_S^{1/3} N_R^{-0.081} N_{Pe}^{-0.715} N_{vdW}^{0.052}
}
\end{equation}

\paragraph{2. Mechanical Interception Retention}
Pure mechanical interception effects:
\begin{equation}
\boxed{
\eta_I = 0.55 A_S N_R^{1.675} N_A^{0.125}
}
\end{equation}

\paragraph{3. Gravitational Retention}
Buoyancy-induced trajectory changes affecting particle-collector collisions:
\begin{equation}
\boxed{
\eta_G = 0.22 N_R^{-0.24} N_G^{1.11} N_{vdW}^{0.053}
}
\end{equation}

\subsection{Strain-induced Blockage Component}

Physical straining effects due to pore throat restrictions:

\begin{equation}
\boxed{
\lambda_{strain} = \frac{k_{str}\psi_{str}}{|w_D + w_{ss}|}
}
\end{equation}

where:
\begin{align}
k_{str} &= 269.7\left(\frac{d_{MP}}{d_g}\right)^{1.42} \\[0.5em]
\psi_{str} &= \left(\frac{d_{MP}+z}{d_{MP}}\right)^{-0.43}
\end{align}

with $z$ measured from sediment-water interface to the layer considered.

\section{Dimensionless Parameters}

\subsection{Porosity-dependent Parameter}

\begin{align}
A_S &= \frac{2(1-\gamma^5)}{2-3\gamma+3\gamma^5-2\gamma^6} \\[0.5em]
\gamma &= (1-\phi)^{1/3}
\end{align}

\subsection{Size and Flow Parameters}

\paragraph{Size Ratio:}
\begin{equation}
N_R = \frac{d_{MP}}{d_g}
\end{equation}

\paragraph{Peclet Number (convection vs diffusion):}
\begin{equation}
N_{Pe} = \frac{|w_D|d_g}{D_L}
\end{equation}

\paragraph{Gravity Number:}
\begin{equation}
N_G = \frac{|w_{ss}|}{|w_D|}
\end{equation}

\subsection{Interaction Parameters}

\paragraph{Van der Waals Number:}
\begin{equation}
N_{vdW} = \frac{A}{kT}
\end{equation}

\paragraph{Attraction Number:}
\begin{equation}
N_A = \frac{A}{12\pi\mu(d_{MP}/2)^2|w_D|}
\end{equation}

\subsection{Parameter Definitions}

\begin{itemize}
    \item $\phi$: Porosity [dimensionless]
    \item $d_{MP}$: Microplastic particle diameter [m]
    \item $A$: Hamaker constant [J]
    \item $k$: Boltzmann constant = $1.381 \times 10^{-23}$ [J/K]
    \item $T$: Temperature [K]
    \item $\mu$: Dynamic viscosity of water [Pa·s]
\end{itemize}

\section{Implementation Notes}

\subsection{Coordinate System}
\begin{itemize}
    \item Vertical coordinate $z$ implemented using FVCOM's sigma coordinate system
    \item Positive $z$ direction: upward from sediment-water interface
    \item Seepage velocity $w_D$ can be positive (upward) or negative (downward)
\end{itemize}

\subsection{Boundary Conditions}

The model supports two distinct simulation modes with different boundary condition treatments for various experimental and natural scenarios.

\subsubsection{Mode 1: Natural Penetration Mode}

This mode simulates natural infiltration of microplastics and toxics from surface water into the sediment column.

\paragraph{Upper Boundary Conditions (z = 0, sediment-water interface)}

\textbf{Movable MP and Embedded Toxics} - Robin boundary condition:
\begin{align}
w_D c_{MP}|_{bot} &= \left[-D_L \frac{\partial p_{MP}}{\partial z} + w_D p_{MP}\right]\bigg|_{z=0} \\[0.5em]
w_D c_T^e|_{bot} &= \left[-D_L \frac{\partial p_T^e}{\partial z} + w_D p_T^e\right]\bigg|_{z=0}
\end{align}

\textbf{Free Toxics} - Dirichlet boundary condition:
\begin{equation}
c_T^f\big|_{z=0} = c_T^f|_{bot}
\end{equation}

\textbf{Blocked MP and Embedded Toxics} - Neumann boundary condition:
\begin{align}
\frac{\partial s_{MP}}{\partial z}\bigg|_{z=0} &= 0 \\[0.5em]
\frac{\partial s_T^e}{\partial z}\bigg|_{z=0} &= 0
\end{align}

\paragraph{Lower Boundary Conditions (z = D_s, bottom outlet)}

\textbf{All tracers} - Neumann boundary condition:
\begin{equation}
\frac{\partial \phi}{\partial z}\bigg|_{z=D_s} = 0 \quad \text{for all tracers } \phi
\end{equation}

\subsubsection{Mode 2: Laboratory Leaching Experiment Mode}

This mode simulates controlled laboratory experiments with water inflow and toxic outflow, with microplastics pre-distributed in the column.

\paragraph{Upper Boundary Conditions (z = 0, inlet)}

\textbf{Free Toxics} - Dirichlet boundary condition:
\begin{equation}
c_T^f\big|_{z=0} = 0
\end{equation}

\textbf{All other tracers} - Neumann boundary condition:
\begin{equation}
\frac{\partial \phi}{\partial z}\bigg|_{z=0} = 0 \quad \text{for } p_{MP}, s_{MP}, p_T^e, s_T^e
\end{equation}

\paragraph{Lower Boundary Conditions (z = D_s, outlet)}

\textbf{Movable MP and Embedded Toxics} - Robin boundary condition (no mass flux):
\begin{align}
\left[-D_L \frac{\partial p_{MP}}{\partial z} + w_D p_{MP}\right]\bigg|_{z=D_s} &= 0 \\[0.5em]
\left[-D_L \frac{\partial p_T^e}{\partial z} + w_D p_T^e\right]\bigg|_{z=D_s} &= 0
\end{align}

\textbf{Free Toxics} - Robin boundary condition (free outflow):
\begin{equation}
w_D c_T^f|_{z=D_s} = \left[-D_L \frac{\partial c_T^f}{\partial z} + w_D c_T^f\right]\bigg|_{z=D_s}
\end{equation}

\textbf{Blocked MP and Embedded Toxics} - Neumann boundary condition:
\begin{align}
\frac{\partial s_{MP}}{\partial z}\bigg|_{z=D_s} &= 0 \\[0.5em]
\frac{\partial s_T^e}{\partial z}\bigg|_{z=D_s} &= 0
\end{align}

\subsection{Boundary Condition Variable Definitions}

\begin{itemize}
    \item $c_{MP}|_{bot}$: MP concentration in bottom layer of surface water [mass/volume]
    \item $c_T^e|_{bot}$: Embedded toxic concentration in bottom layer of surface water [mass/volume]
    \item $c_T^f|_{bot}$: Free toxic concentration in bottom layer of surface water [mass/volume]
    \item $z = 0$: Sediment-water interface (uppermost sediment layer) [m]
    \item $D_s$: Total thickness of sediment matrix layer [m]
\end{itemize}

\subsection{Boundary Condition Types}
\begin{itemize}
    \item \textbf{Robin conditions}: Relate advective-diffusive flux to concentration gradients
    \item \textbf{Dirichlet conditions}: Specify concentration values directly
    \item \textbf{Neumann conditions}: Specify zero gradient (no flux through boundary)
\end{itemize}

\subsection{Initial Conditions Framework}

The model supports multiple initialization approaches for sediment column concentrations.

\subsubsection{Uniform Concentration Initialization}

Read uniform concentrations from \texttt{generic\_plastic.inp} parameters:

\begin{verbatim}
! Sediment column initial concentrations (uniform)
INIT_MP_CONC_SEDIMENT = 0.0001        ! kg/m^3
INIT_ETOX_CONC_SEDIMENT = 0.00001     ! kg/m^3  
INIT_FTOX_CONC_SEDIMENT = 0.0         ! kg/m^3
INIT_SMASS_DISTRIBUTION = 0.5         ! fraction in blocked phase
\end{verbatim}

\subsubsection{NetCDF File Initialization}

Read spatially-varying initial conditions from NetCDF files via namelist:

\begin{verbatim}
&NML_ADDITIONAL_MODELS
INIT_INFILTRATION_TYPE = 'netcdf'
INIT_INFILTRATION_FILE = 'initial_sediment_profiles.nc'
\end{verbatim}

\paragraph{NetCDF File Structure:}
\begin{itemize}
    \item \texttt{mp\_sediment\_conc(node, layer, nplast)}: Movable MP concentrations
    \item \texttt{mp\_sediment\_blocked(node, layer, nplast)}: Blocked MP concentrations  
    \item \texttt{etox\_sediment\_conc(node, layer, nplast)}: Embedded toxic concentrations
    \item \texttt{etox\_sediment\_blocked(node, layer, nplast)}: Blocked embedded toxics
    \item \texttt{ftox\_sediment\_conc(node, layer)}: Free toxic concentrations
\end{itemize}

\subsubsection{Implementation in init\_plast.F}

\paragraph{Current Structure (lines 69-97):}
\begin{itemize}
    \item \texttt{LOAD\_INIT\_PLAST\_MASS()}: Bottom mass initialization
    \item \texttt{LOAD\_INIT\_TOXIN\_MASS()}: Toxic mass initialization
\end{itemize}

\paragraph{Extensions Needed:}
\begin{itemize}
    \item \texttt{LOAD\_INIT\_SEDIMENT\_PROFILES()}: 3D concentration profiles
    \item Layer-by-layer initialization for each tracer type
    \item Boundary condition coupling with water column
\end{itemize}

\section{Embedded Toxic Transport Equations}

Embedded toxics are carried by microplastics and follow identical transport dynamics with the same physical parameters and leaching kinetics.

\subsection{Primary Embedded Toxic Transport Equation}

The transport equation for movable embedded toxics ($p_T^e$) has identical form to the MP equation:

\begin{equation}
\boxed{
\frac{\partial p_T^e}{\partial t} + (w_D + w_{ss}) \frac{\partial p_T^e}{\partial z} = \frac{\partial}{\partial z}\left[D_L \frac{\partial p_T^e}{\partial z}\right] - \lambda|w_D + w_{ss}|p_T^e - S_{leach}^p
}
\end{equation}

\subsection{Blocked Embedded Toxic Equation}

The evolution equation for blocked embedded toxics ($s_T^e$):

\begin{equation}
\boxed{
\frac{\partial s_T^e}{\partial t} = \lambda|w_D + w_{ss}|p_T^e - S_{leach}^s
}
\end{equation}

\subsection{Total Embedded Toxic Concentration}

The total embedded toxic concentration in the sediment is:

\begin{equation}
\boxed{
c_T^e = p_T^e + s_T^e
}
\end{equation}

\subsection{Variable Definitions}

\begin{itemize}
    \item $p_T^e$: Movable embedded toxic concentration in sediment [mass/volume]
    \item $s_T^e$: Blocked embedded toxic concentration (immobilized by solid matrix) [mass/volume]
    \item $c_T^e$: Total embedded toxic concentration in sediment [mass/volume]
\end{itemize}

\subsection{Key Features}

\begin{itemize}
    \item \textbf{Identical Transport Parameters}: All parameters ($w_D$, $w_{ss}$, $D_L$, $\lambda$) are identical to MP transport since toxics are carried by MPs
    \item \textbf{Same Leaching Kinetics}: Leaching terms ($S_{leach}^p$, $S_{leach}^s$) follow the same formulation as MP system
    \item \textbf{Same Retention Mechanisms}: DLVO theory and strain-induced blocking apply identically
    \item \textbf{Coupled Transport}: Physical coupling with MP transport through identical governing physics
\end{itemize}

\section{Free Toxic Transport Equations}

Free toxics are not subject to blocking mechanisms and follow simplified transport dynamics driven by seepage flow and leaching from microplastics.

\subsection{Free Toxic Transport Equation}

The governing equation for free toxics ($c_T^f$) incorporates advection, diffusion, and leaching source terms:

\begin{equation}
\boxed{
\frac{\partial c_T^f}{\partial t} + w_D \frac{\partial c_T^f}{\partial z} = \frac{\partial}{\partial z}\left[D_L \frac{\partial c_T^f}{\partial z}\right] + S_{leach}
}
\end{equation}

\subsection{Leaching Source Term}

The leaching source combines toxin release from both blocked and movable MP phases:

\begin{equation}
\boxed{
S_{leach} = S_{leach}^p + S_{leach}^s
}
\end{equation}

where:
\begin{itemize}
    \item $S_{leach}^p$: Leaching from movable MP phase [mass/(volume·time)]
    \item $S_{leach}^s$: Leaching from blocked MP phase [mass/(volume·time)]
\end{itemize}

\subsection{Variable Definitions}

\begin{itemize}
    \item $c_T^f$: Mass concentration of free toxics in sediment [mass/volume]
    \item $w_D$: Darcy velocity (seepage flow, identical to MP/embedded toxic equations) [m/s]
    \item $D_L$: Mechanical seepage diffusivity (identical to MP/embedded toxic equations) [m$^2$/s]
    \item $S_{leach}$: Combined leaching source from blocked and movable phases [mass/(volume·time)]
\end{itemize}

\subsection{Key Features}

\begin{itemize}
    \item \textbf{No Settling Velocity}: $w_{ss} = 0$ - free toxics don't settle independently
    \item \textbf{No Blocking Mechanisms}: Free toxics are not retained by DLVO or strain effects
    \item \textbf{Pure Advection-Diffusion}: Transport driven by seepage flow and diffusion only
    \item \textbf{Source-Driven Dynamics}: Concentration evolution controlled by MP leaching processes
    \item \textbf{Same Physical Diffusivity}: Uses identical $D_L$ formulation as other transport phases
    \item \textbf{Simplified Transport}: No retention parameters ($\lambda$) or complex retention physics
\end{itemize}

\section{Numerical Implementation Strategy}

\subsection{Existing Numerical Methods (Reference: mod\_plast.F)}

The infiltration module leverages existing proven numerical schemes from the FVCOM-MP framework.

\subsubsection{Flux-Limited Settling Algorithm}

\paragraph{Reference:} \texttt{Settle\_Flux\_new()} subroutine (mod\_plast.F:1829-1982)

\paragraph{Key Features:}
\begin{itemize}
    \item \textbf{CFL-limited time stepping}: $\Delta t_{max} = CFL_{settle} \cdot \min(\Delta x_i/w_{set,i})$
    \item \textbf{Multi-cycle approach}: $n_{cycle} = \lfloor \Delta t_{plast}/\Delta t_{max} + 1 \rfloor$
    \item \textbf{Flux limiters}: Support for upwind, minmod, Van Leer limiters
    \item \textbf{Bottom boundary handling}: Different treatments for boundary conditions
\end{itemize}

\paragraph{Adaptation for Infiltration:}
\begin{itemize}
    \item Replace settling velocity $w_{set}$ with combined seepage-settling: $w_D + w_{ss}$
    \item Modify boundary conditions for sediment-water interface
    \item Add blocking term integration: $-\lambda|w_D + w_{ss}|p_{MP}$
\end{itemize}

\subsubsection{Vertical Diffusion Treatment}

\paragraph{Reference:} \texttt{vdif\_scal()} calls (mod\_plast.F:863, 907)

\paragraph{Key Features:}
\begin{itemize}
    \item \textbf{Implicit scheme}: Handles stiff diffusion problems
    \item \textbf{Tridiagonal solver}: Efficient 1D solution
    \item \textbf{Multiple tracers}: Supports different diffusivity values
\end{itemize}

\paragraph{Adaptation for Infiltration:}
\begin{itemize}
    \item Use mechanical diffusivity $D_L$ instead of turbulent diffusivity
    \item Apply to all infiltration tracers: $p_{MP}$, $s_{MP}$, $p_T^e$, $s_T^e$, $c_T^f$
\end{itemize}

\subsubsection{Time Integration Strategy}

\paragraph{Current Approach (mod\_plast.F:749-950):}
\begin{itemize}
    \item \textbf{Operator splitting}: Advection $\rightarrow$ Diffusion $\rightarrow$ Reactions
    \item \textbf{Multi-step updates}: \texttt{cnew} $\rightarrow$ \texttt{conc} progression
    \item \textbf{Source term integration}: Explicit treatment
\end{itemize}

\paragraph{Infiltration Integration (pseudo-code):}
\begin{verbatim}
! Infiltration time-stepping
call calc_infiltration_advection(p_MP, w_D + w_ss, DT_infil)
call calc_infiltration_diffusion(p_MP, D_L, DT_infil)  
call calc_blocking_reactions(p_MP, s_MP, lambda, DT_infil)
call calc_toxic_leaching(p_T_e, s_T_e, c_T_f, S_leach, DT_infil)
\end{verbatim}

\subsection{Proposed Infiltration Subroutines}

\subsubsection{Main Integration: \texttt{Calc\_Infiltration()}}
\begin{itemize}
    \item \textbf{Location}: Add after \texttt{Calc\_Deposition} in mod\_plast.F
    \item \textbf{Function}: Coordinate all infiltration processes
    \item \textbf{Structure}: Loop over nodes and plastic classes, apply boundary conditions, call process subroutines, update arrays
\end{itemize}

\subsubsection{Transport Solver: \texttt{Infiltrate\_1D\_Transport()}}
\begin{itemize}
    \item \textbf{Function}: Solve 1D advection-diffusion-reaction system
    \item \textbf{Method}: Based on existing \texttt{Settle\_Flux\_new()} framework
    \item \textbf{Features}: CFL-limited time stepping, flux limiters, Robin/Dirichlet/Neumann boundary conditions
\end{itemize}

\subsubsection{Blocking Kinetics: \texttt{Calc\_Retention\_Kinetics()}}
\begin{itemize}
    \item \textbf{Function}: Handle DLVO + strain blocking mechanisms
    \item \textbf{Method}: Explicit integration of blocking rates
    \item \textbf{Components}: Calculate dimensionless parameters, compute retention efficiencies, apply blocking/release rates
\end{itemize}

\subsubsection{Boundary Coupling: \texttt{Apply\_Infiltration\_BC()}}
\begin{itemize}
    \item \textbf{Function}: Interface with water column concentrations
    \item \textbf{Features}: Mode selection (natural vs laboratory), Robin condition implementation, mass balance enforcement
\end{itemize}

\subsection{Discretization Approach}
Based on FVCOM finite volume framework with sigma coordinate system for vertical discretization.

\subsection{Coupling with Water Column Transport}
Interface conditions implemented through boundary flux coupling between water column settling and sediment infiltration processes.

\section{Model Parameters Summary}

\begin{table}[h!]
\centering
\begin{tabular}{@{}lll@{}}
\toprule
\textbf{Parameter} & \textbf{Description} & \textbf{Units} \\
\midrule
\multicolumn{3}{c}{\textbf{Transport Parameters}} \\
\midrule
$w_D$ & Darcy velocity (seepage) & m/s \\
$w_{ss}$ & MP settling velocity in sediment & m/s \\
$D_L$ & Mechanical seepage diffusivity & m$^2$/s \\
$\lambda$ & Total blockage rate & 1/s \\
\midrule
\multicolumn{3}{c}{\textbf{Concentration Variables}} \\
\midrule
$p_{MP}$ & Movable MP concentration & mass/volume \\
$s_{MP}$ & Blocked MP concentration & mass/volume \\
$c_{MP}$ & Total MP concentration & mass/volume \\
$p_T^e$ & Movable embedded toxic concentration & mass/volume \\
$s_T^e$ & Blocked embedded toxic concentration & mass/volume \\
$c_T^e$ & Total embedded toxic concentration & mass/volume \\
$c_T^f$ & Free toxic concentration & mass/volume \\
\midrule
\multicolumn{3}{c}{\textbf{Physical Parameters}} \\
\midrule
$\phi$ & Sediment porosity & - \\
$d_{MP}$ & Microplastic diameter & m \\
$d_g$ & Sediment grain size & m \\
$A$ & Hamaker constant & J \\
$\mu$ & Water dynamic viscosity & Pa·s \\
$T$ & Temperature & K \\
$c$ & Mechanical diffusivity coefficient & - \\
$\tau$ & Tortuosity & - \\
$D_m$ & Molecular diffusivity & m$^2$/s \\
\bottomrule
\end{tabular}
\caption{Key model parameters, variables, and their units}
\label{tab:parameters}
\end{table}

\section{References}

\textcolor{red}{[To be added based on literature cited in development]}

\appendix

\section{Derivation Details}

\textcolor{red}{[Detailed derivations of key equations if needed]}

\section{Validation Test Cases}

\textcolor{red}{[Analytical solutions for model validation]}

\end{document}